\section{Verlet Integration \& Core Kinematics}

This section documents the foundational Verlet integration method used for particle position updates, along with derived kinematic quantities such as velocity, acceleration, and angular motion.

\cite{jakobsen2001advanced} presents Verlet integration as a simple yet powerful technique for physics simulation, encoding velocity implicitly in position differences rather than explicitly tracking velocity. This section derives all kinematic equations from the core Verlet update.

\subsection{Verlet Position Integration}

\begin{equation}
  \vv{x}_{\text{new}} = \vv{x}_{\text{current}} + (\vv{x}_{\text{current}} - \vv{x}_{\text{prev}}) + \vv{a} \cdot dt^2
  \label{eq:verlet-position}
\end{equation}

\textbf{Physical Meaning}: The new position is computed from the current position, the velocity (implicit in the difference $\vv{x}_{\text{current}} - \vv{x}_{\text{prev}}$), and the acceleration. This formulation is stable and does not require explicit velocity tracking.

\textbf{Source Code}: \coderef{physics.js}{1000--1018}

\begin{lstlisting}[caption={Verlet position update with rotational displacement}]
function verletIntegrateAllPoints(dt, netAccelX, netAccelY, netAngularAccel) {
  for (let i = 0; i < 4; i++) {
    const p = particles[i];
    const dx = p.x - p.xPrev;
    const dy = p.y - p.yPrev;
    const armX = p.x - body.center.x;
    const armY = p.y - body.center.y;

    p.xPrev = p.x;
    p.yPrev = p.y;

    p.x += dx + (netAccelX + (-netAngularAccel * armY)) * dt * dt;
    p.y += dy + (netAccelY + (netAngularAccel * armX)) * dt * dt;
  }
}
\end{lstlisting}

\textbf{Parameters}:
\begin{itemize}
  \item $\vv{x}_{\text{current}}$: Current particle position $\begin{pmatrix} x \\ y \end{pmatrix}$
  \item $\vv{x}_{\text{prev}}$: Previous particle position
  \item $\vv{a}$: Linear acceleration $\begin{pmatrix} a_x \\ a_y \end{pmatrix}$ (includes all forces)
  \item $dt$: Fixed physics timestep (typically 0.01 s for 100 Hz)
\end{itemize}

\textbf{Notes}:
\begin{itemize}
  \item The formulation is \textit{symplectic} (energy-conserving) for conservative forces
  \item Velocity is implicit: $\vv{v} = (\vv{x}_{\text{current}} - \vv{x}_{\text{prev}}) / dt$
  \item Rotation is handled by applying angular acceleration to all particles relative to the center of mass
  \item The constraint solver (Section \ref{sec:constraints}) post-processes positions to maintain rigid-body shape
\end{itemize}

\subsection{Velocity Extraction}

\begin{equation}
  \vv{v} = \frac{\vv{x}_{\text{current}} - \vv{x}_{\text{prev}}}{dt}
  \label{eq:velocity-extraction}
\end{equation}

\textbf{Physical Meaning}: Velocity is computed backward-in-time from the position difference. This provides a natural approximation of instantaneous velocity.

\textbf{Source Code}: \coderef{physics.js}{155--165}

\textbf{Notes}:
\begin{itemize}
  \item Velocity is used for tire slip angle calculation, drag force, and sensor feedback (gauges)
  \item Per-wheel velocity: $\vv{v}_{\text{wheel}} = \vv{v}_{\text{body}} + \omega \times (\vv{p}_{\text{wheel}} - \vv{c}_{\text{cm}})$
\end{itemize}

\subsection{Angular Velocity}

\begin{equation}
  \omega = \frac{\Delta \theta}{dt}
  \label{eq:angular-velocity}
\end{equation}

\textbf{Physical Meaning}: Angular velocity is the rate of change of heading angle. It is derived from consecutive heading values wrapped to $(-\pi, \pi]$.

\textbf{Source Code}: \coderef{physics.js}{180--190}

\begin{lstlisting}[caption={Angular velocity computation with angle wrapping}]
const headingDelta = wrapAngle(heading - prevHeading);
angularVelocity = headingDelta / dt;
\end{lstlisting}

\textbf{Notes}:
\begin{itemize}
  \item Heading is computed from wheel positions: $\theta = \text{atan2}(\text{front}_y - \text{rear}_y, \text{front}_x - \text{rear}_x)$
  \item Angular velocity is used for yaw damping, steering feedback, and motion blur effects
\end{itemize}

\subsection{Linear Acceleration}

\begin{equation}
  \vv{a} = \frac{\vv{v}_{\text{current}} - \vv{v}_{\text{prev}}}{dt}
  \label{eq:linear-acceleration}
\end{equation}

\textbf{Physical Meaning}: Linear acceleration is the rate of change of velocity, computed from consecutive velocity snapshots.

\textbf{Source Code}: \coderef{physics.js}{200--210}

\textbf{Notes}:
\begin{itemize}
  \item Acceleration is typically filtered via exponential moving average (EMA) to reduce noise before display
  \item Longitudinal acceleration: $a_{\text{long}} = \vv{a} \cdot \hat{\vv{u}}_{\text{forward}}$
  \item Lateral acceleration: $a_{\text{lat}} = \vv{a} \cdot \hat{\vv{u}}_{\text{right}}$
  \item Used for weight transfer calculations and feedback displays
\end{itemize}

\subsection{Jerk (Third Derivative)}

\begin{equation}
  j = \frac{a_{\text{current}} - a_{\text{prev}}}{dt}
  \label{eq:jerk}
\end{equation}

\textbf{Physical Meaning}: Jerk is the rate of change of acceleration, representing the smoothness of motion transitions. High jerk indicates sudden changes in acceleration.

\textbf{Source Code}: \coderef{physics.js}{215--225}

\textbf{Notes}:
\begin{itemize}
  \item v14 fixed a critical bug in jerk computation: v9 incorrectly divided by $dt$ (100× too large)
  \item v14 correctly computes: $j = (a - a_{\text{prev}}) / dt$
  \item Jerk is used for skid mark visual intensity and motion blur effects
  \item Filtered jerk is displayed in diagnostic HUD
\end{itemize}

\newpage
